\section{Dynamical Symmetries and Conservation Laws}

 [... non lie symmetry content ...]

When a dynamical symmetry group \(G\) is a Lie group, with a parametrization \(g(\underline{\theta})\), the quantum unitary representation \(\tilde{U}(\underline{\theta})\) representing the action of \(G\) on the Hilbert space \(\mathcal{H}\) respects
\[
    \tilde{U}(\underline{\theta})  H \tilde{U}(\underline{\theta})^\dagger = H, \quad \forall \underline{\theta}.
\]
This property restricts the form of the Hamiltonian \(H\) on one hand and has consequences for the infinitesimal generators \(Q_a\) of the unitary family \(\tilde{U}(\underline{\theta})\) on the other: \(H\) commutes with the \(Q_a\)
\[
    [H, Q_a] = 0.
\]
\begin{proof}
    If we compute the derivative of the invariance condition with respect to \(\theta^a\) at \(\underline{\theta} = 0\), we get
    \[
        \begin{aligned}
            0 & = i \frac{\partial}{\partial \theta^a} (H - \tilde{U}(\underline{\theta}) H \tilde{U}(\underline{\theta})^\dagger) \bigg|_{\underline{\theta} = \underline{0}} = \left[ -i \frac{\partial \tilde{U}(\underline{\theta})}{\partial \theta^a} H \tilde{U}(\underline{\theta})^\dagger + \tilde{U}(\underline{\theta})H\left( i \frac{\partial \tilde{U}(\underline{\theta})}{\partial \theta^a}\right)^{\dagger} \right]_{\underline{\theta} = \underline{0}} \\
              & = -i \frac{\partial \tilde{U}(\underline{0})}{\partial \theta^a} H \tilde{U}(\underline{0})^\dagger + \tilde{U}(\underline{0})H\left( i \frac{\partial \tilde{U}(\underline{0})}{\partial \theta^a}\right)^{\dagger} = - Q_a H + H Q_a = [H, Q_a],
        \end{aligned}
    \]
    since \(\tilde{U}(\underline{0}) = \mathbb{I}\) and \(Q_a = -i \frac{\partial \tilde{U}(\underline{\theta})}{\partial \theta^a} \big|_{\underline{\theta} = \underline{0}}\). This proves that when a system is invariant under a Lie group of dynamical symmetries, the generators of that group commute with the Hamiltonian.
\end{proof}

This result is significant because it implies that the generators \(Q_a\) are conserved quantities of the system, as they commute with the Hamiltonian. They are obserables with meaningful physical significance. In quantum mechanics, this means that the eigenvalues of \(Q_a\) are \textbf{constants of motion}, and the system's dynamics can be analyzed in terms of these conserved quantities. This connection between symmetries and conservation laws is a fundamental aspect of both classical and quantum physics, encapsulated in \textbf{Noether's theorem}. Indeed, we can define the evolution operator
\[
    U(t,s) = e^{-\tfrac{i}{\hbar}(t-s)H},
\]
such that any observable \(\hat{A}\) evolves in time according to
\[
    \hat{A}(t) = U(t,0)^\dagger \hat{A} U(t,0).
\]
Then we have
\[
    U(t,s) Q_a U(t,s)^\dagger = e^{-\tfrac{i}{\hbar}(t-s)H} Q_a e^{\tfrac{i}{\hbar}(t-s)H} = Q_a,
\]
which shows that \(Q_a(t)=Q_a\) is conserved under the time evolution generated by \(H\).
\begin{proof}
    This is a direct consequence of the commutation relation \([H, Q_a] = 0\): since \(Q_a\) commutes with \(H\), it also commutes with any function of \(H\), including the exponential \(e^{-\tfrac{i}{\hbar}(t-s)H}\). Therefore, we can move \(Q_a\) through the exponential without changing its value, leading to the conclusion that \(Q_a\) is invariant under the time evolution generated by \(H\).
\end{proof}
We have just shown the quantum analog of Noether's theorem: for every continuous symmetry of the action (or equivalently, for every Lie group of dynamical symmetries), there is a corresponding conserved quantity. This profound connection between symmetries and conservation laws is a cornerstone of modern physics, providing deep insights into the fundamental structure of physical theories.

The symmetry Lie algebra \(\mathfrak{g}\)
\[
    [t_a, t_b] = f^c_{\ ab} t_c,
\]
always contains a Cartan subalgebra \(\mathfrak{h}\), that is a maximal subspace of \(\mathfrak{g}\) the Lie brackets of all whose elements vanish: all the structure constants \(f^c_{\ ab}\) vanish. Denote by \(h_i\) a set of generators of \(\mathfrak{h}\). These are certain linearly independent linear combinations of the generators \(T_a\) of \(\mathfrak{g}\)
\[
    h_i = h_i^a t_a,
\]
with the property that
\[
    [h_i, h_j] = 0.
\]

Thus we are lead to introduce a set of quantum generators \(C_i\) of the Cartan subalgebra \(\mathfrak{h}\) in such a way that they commute with the Hamiltonian \(H\) and with each other
\[
    C_i = h_i^a Q_a.
\]
These generators \(C_i\), defined in this way, respects
\[
    [C_i, C_j] = 0.
\]

\begin{proof}
    \[
        0 = [h_i, h_j] = [h_i^a t_a, h_j^b t_b]\dots
    \]
    thus we can write
    \[
        \begin{aligned}
            [C_i, C_j] & = [h_i^a Q_a, h_j^b Q_b] = h_i^a \, h_j^b [Q_a, Q_b] \\
                       & = i h_i^a \, h_j^b f^c_{\ ab} Q_c = 0,
        \end{aligned}
    \]
    since the structure constants \(f^c_{\ ab}\) vanish for all \(a, b\) when \(t_a\) and \(t_b\) are in the Cartan subalgebra \(\mathfrak{h}\).
\end{proof}

Thus
\[
    [C_i, C_j] = 0, \quad [C_i, H] = 0.
\]
There is a one-to-one correspondence between the eigenvalues of the \(C_i\) and the generators of the cartan subalgebra \(\mathfrak{h}\).

Consider now
\[
    \eta_{ab} = \text{Tr}(t_a t_b),
\]
if \(\mathfrak{g}\) is a semisimple Lie algebra, the matrix \(\eta_{ab}\) is non-degenerate and can be used to raise and lower indices. It will be invertible and we can define \(\eta^{ab}\) such that \(\eta^{ab} \eta_{bc} = \delta^a_c\). We obtained a metric on the Lie algebra \(\mathfrak{g}\) and we can use.

Now we are going to prove that the structure constants \(f_{abc} = f^d_{\ bc} \eta_{da}\) of a semisimple Lie algebra are completely antisymmetric in the indices \(a, b, c\) and not only the last two.
\begin{proof}
    Since
    \[
        \Tr([A, B]) = \Tr(AB - BA) = 0,
    \]
    we can write
    \[
        \begin{aligned}
            0 & = \Tr([t_a, t_b t_c]) = \Tr([t_a, t_b] t_c + t_b [t_a, t_c]) = \Tr(f^d_{\ ab} t_d t_c + f^d_{\ ac} t_b t_d) \\
              & = f^d_{\ ab} \eta_{dc} + f^d_{\ ac} \eta_{db} = f_{abc} + f_{acb} = 0,
        \end{aligned}
    \]
    or equivalently
    \[
        f^{a \ c}_{\ b} + f^{b \ a}_{\ c} = 0.
    \]
\end{proof}

The antisymmetry of the structure constants \(f_{abc}\) is a crucial property that has significant implications for the representation theory of Lie algebras and the classification of semisimple Lie algebras. \(\eta_{ab}\) is called the \textbf{Cartan metric} of the Lie algebra \(\mathfrak{g}\) and through it we can define the Quadratic Casimir operator \(C_2\) of the Lie algebra \(\mathfrak{g}\)
\[
    C_2 = \eta^{ab} Q_a Q_b.
\]
The Casimir operator \(C_2\) is a central element of the universal enveloping algebra of \(\mathfrak{g}\), meaning that it commutes with all elements of \(\mathfrak{g}\). In particular, since \(Q_a\) are the generators of the symmetry group, we have
\[
    \begin{aligned}
        [Q_a, C_2] = [Q_a, \eta^{bc} Q_b Q_c] = \eta^{bc} ([Q_a, Q_b] Q_c + Q_b [Q_a, Q_c])                               \\
        = i \eta^{bc} (f^d_{\ ab} Q_d Q_c + f^d_{\ ac} Q_b Q_d) = i ( f^{d \ b}_{\ a} Q_d Q_b + f^{d \ b}_{\ a} ) Q_b Q_d \\
        = i f^{d \ b}_{\ a} (Q_d Q_b + Q_b Q_d) = 0,
    \end{aligned}
\]
where we have used the antisymmetry of the structure constants \(f_{abc}\) in the last step. The Casimir operator \(C_2\) is an important quantity in the representation theory of Lie algebras, as it takes a constant value on each irreducible representation, providing a way to classify representations and understand their properties.
\[
    [C_i, C_j] = 0, \quad [H, C_i]=0, \quad [C_2, C_i] = 0, \quad [H, C_2] = 0.
\]

To sum up, we have a dynmical lie symmetry group \(G\) with Lie algebra \(\mathfrak{g}\) and a Cartan subalgebra \(\mathfrak{h}\) of \(\mathfrak{g}\). We have a set of generators \(C_i\) of the Cartan subalgebra \(\mathfrak{h}\) that
\[
    \begin{dcases}
        C_i = h_i^a Q_a, \\
        C_2 = \eta^{ab} Q_a Q_b,
    \end{dcases}
\]
commute with the Hamiltonian \(H\) and with each other, and we have the Casimir operator \(C_2\) that also commutes with \(H\) and with the \(C_i\). The eigenvalues of the \(C_i\) and \(C_2\) can be used to label the states of the system and to classify the irreducible representations of the symmetry group \(G\).

\begin{example}[The Hydrogen Atom]
    The hydrogen atom has a dynamical symmetry group known as the \textbf{SO(3)} group, associated to the algebra \(\mathfrak{so}(3)\) with generators
    \[
        [t_a, t_b] = \epsilon_{abc} t_c.
    \]
    We have an associated cartan subalgebra generated by \(c t_3\), with \(c \in \mathbb{R}\)  and a cartesian metric \(\eta_{ab} = \Tr(t_a t_b) = \delta_{ab}\). The quantum generators \(J_i\) (analogue of \(Q_a\)) of the symmetry group commute with the Hamiltonian of the hydrogen atom, and the analogue of \(C_i\) is \(J_3\), while the analogue of the Casimir operator \(C_2\) is \(J^2 = J_1^2 + J_2^2 + J_3^2 = \delta_{ij} J_i J_j\). The eigenvalues of \(J_3\) and \(J^2\) can be used to label the states of the hydrogen atom and to classify the irreducible representations of the SO(3) symmetry group, which correspond to the different energy levels and angular momentum states of the electron in the hydrogen atom.

    Thus
    \[
        \begin{aligned}
            H \Phi_{w c \underline{e}} = \Phi_{w c \underline{e}} w,   \\
            C_2 \Phi_{w c \underline{e}} = \Phi_{w c \underline{e}} c, \\
            C_i \Phi_{w c \underline{e}} = \Phi_{w c \underline{e}} e_i,
        \end{aligned}
    \]
    where \(\Phi_{w c \underline{e}}\) are the eigenstates of the Hamiltonian \(H\), the Casimir operator \(C_2\) and the generators \(C_i\) of the Cartan subalgebra, with eigenvalues \(w\), \(c\) and \(e_i\) respectively. The quantum numbers \(w\), \(c\) and \(e_i\) can be used to classify the states of the hydrogen atom and to understand their properties in terms of the underlying symmetries of the system.

    Using \(w\) as the Bohr energy levels, we can write the energy eigenvalues of the hydrogen atom as
    \[
        \begin{aligned}
            w   & \to n, \quad n = 1, 2, 3, \dots,          \\
            c   & \to l, \quad l = 0, 1, 2, \dots, n-1,     \\
            e_i & \to m, \quad m = -l, -l+1, \dots, l-1, l.
        \end{aligned}
    \]
    where in this case \(\underline{e} = e\) is a single quantum number since we have only one generator \(C_i\) of the Cartan subalgebra. The energy levels of the hydrogen atom are quantized and depend on the principal quantum number \(n\), the angular momentum quantum number \(l\) and the magnetic quantum number \(m\), with the energy eigenvalues given by
    \[
        c = \hbar^2 l (l+1), \quad \dots
    \]
\end{example}

We obtained a bunch of commuting self adjoint operators \(H, C_i, C_2\) that can be simultaneously diagonalized, with common eigenvectors \(\Phi_{w c \underline{e}}\) and eigenvalues \(w, c, e_i\). The quantum numbers \(w, c, e_i\) can be used to label the states of the system and to classify the irreducible representations of the symmetry group \(G\). This is a powerful tool for understanding the structure of quantum systems and for predicting their behavior based on their symmetries.

We found an orthonormal basis of eigenvectors \(\{\Phi_{w c \underline{e}}\}\) of the Hilbert space \(\mathcal{H}\) that are common eigenvectors of the Hamiltonian \(H\), the Casimir operator \(C_2\) and the generators \(C_i\) of the Cartan subalgebra:
\[
    \begin{aligned}
        H \Phi_{w c \underline{e}}   & = w \Phi_{w c \underline{e}},   \\
        C_2 \Phi_{w c \underline{e}} & = c \Phi_{w c \underline{e}},   \\
        C_i \Phi_{w c \underline{e}} & = e_i \Phi_{w c \underline{e}},
    \end{aligned}
\]
This basis allows us to express any state of the system as a linear combination of these eigenvectors, and to analyze the dynamics and properties of the system in terms of its symmetries and conserved quantities.

